\paragraph{各向异性 Serrin 射线的 Fibonacci 同余骨架、局部化账本与六倍窗边界阈值}
在 Figalli--Zhang 对 ADR 与全局 \(\beta\)-square-function 假设下的各向异性 Serrin 刚性已经把可实现粗糙域压缩为 Wulff 形的平移与伸缩，并且折叠同余的模半环表示、window-\(6\) 的 \(18\oplus 3\) 刚性分裂、boundary sheet-flip 的规范中心、容量曲线与纤维直方图的完备统计、局部化数系的 Pontryagin 对偶分解以及六倍窗零点块下界都已建立之后，还可进一步得到如下七码接口结果。它们分别刻画未归一化 Serrin 射线的有限分辨率模半环、该射线上的精确容量律、局部化 Minkowski 射线的圆因子--素数账本分解、window-\(6\) boundary 三点集的算术钉扎、boundary 扇区的 Weyl 完成阈值、六倍窗算术零点块在 Serrin 射线骨架上的传送，以及单参数各向异性 Serrin 观测的 Fibonacci 同余普适性。以下命题均是对前述结果的推导性转写。 \cite{FigalliZhang2026AnisotropicSerrinGeneralDomains}

\begin{theorem}[未归一化各向异性 Serrin 射线的有限分辨率模半环与 Fibonacci 同余完备化]\label{thm:xi-time-part9x-serrin-ray-fibonacci-semirings}
固定满足 Figalli--Zhang 假设的均匀凸各向异性体 \(K\)，记对应 Wulff 形为 \(W_K\)。模去平移并固定重心于原点后，添入形式零元
\[
0\cdot W_K:=\{0\},
\]
并定义整数尺度射线
\[
\mathfrak W_K^{(0)}:=\{\rho W_K:\ \rho\in \NN_0\}.
\]
在其上定义两种运算
\[
(\rho W_K)\boxplus(\sigma W_K):=(\rho+\sigma)W_K,
\qquad
(\rho W_K)\boxtimes(\sigma W_K):=(\rho\sigma)W_K.
\]
对每个 \(m\ge 1\)，再定义同余
\[
\rho W_K\equiv_m^{(K)}\sigma W_K
\iff
\rho\equiv \sigma \pmod{F_{m+2}}.
\]
则映射
\[
[\rho W_K]_{\equiv_m^{(K)}}\longmapsto [\rho]_{F_{m+2}}
\]
给出规范半环同构
\[
\bigl(\mathfrak W_K^{(0)}/\!\equiv_m^{(K)},\boxplus,\boxtimes\bigr)
\cong
\bigl(\ZZ/(F_{m+2}\ZZ),+,\times\bigr).
\]

再记 \(\mathfrak W_K^{\mathrm{gp}}\) 为 \((\mathfrak W_K^{(0)},\boxplus)\) 的 Grothendieck 群完成。则
\[
\mathfrak W_K^{\mathrm{gp}}\cong \ZZ.
\]
若赋予 \(\mathfrak W_K^{\mathrm{gp}}\) 由子群
\[
F_{m+2}\,\mathfrak W_K^{\mathrm{gp}}
\qquad(m\ge 1)
\]
生成的 Fibonacci 同余拓扑，则其完备化满足规范拓扑群同构
\[
\widehat{\mathfrak W}_{K,F}^{\mathrm{gp}}
\cong
\widehat{\ZZ}_F
\cong
\widehat{\ZZ}.
\]
\end{theorem}
\begin{proof}
由 Figalli--Zhang 的主定理，模去平移并固定重心后，一切可实现各向异性 Serrin 域恰为单参数族
\[
\{\rho W_K:\ \rho>0\}.
\]
对同心 Wulff 形，Minkowski 和满足
\[
\rho W_K+\sigma W_K=(\rho+\sigma)W_K,
\]
而尺度复合显然给出
\[
\rho(\sigma W_K)=(\rho\sigma)W_K.
\]
故映射
\[
\iota:\mathfrak W_K^{(0)}\longrightarrow \NN_0,
\qquad
\iota(\rho W_K)=\rho
\]
把 \((\mathfrak W_K^{(0)},\boxplus,\boxtimes)\) 识别为 \((\NN_0,+,\times)\)。

于是对每个 \(m\ge 1\)，同余 \(\equiv_m^{(K)}\) 正是模 \(F_{m+2}\) 同余在 \(\NN_0\) 上的拉回，故
\[
[\rho W_K]_{\equiv_m^{(K)}}\longmapsto [\rho]_{F_{m+2}}
\]
立即给出所述半环同构。

另一方面，\((\NN_0,+)\) 的 Grothendieck 群完成是 \(\ZZ\)，故
\[
\mathfrak W_K^{\mathrm{gp}}\cong \ZZ.
\]
再把由 \(F_{m+2}\mathfrak W_K^{\mathrm{gp}}\) 生成的拓扑经上述同构输运到 \(\ZZ\)，便恰得到定理 \ref{thm:xi-visible-arithmetic-fibonacci-adic-profinite-coincidence} 中的 Fibonacci 同余拓扑 \(\tau_F\)。因此其完备化满足
\[
\widehat{\mathfrak W}_{K,F}^{\mathrm{gp}}
\cong
\widehat{\ZZ}_F
\cong
\widehat{\ZZ}.
\]
证毕。
\end{proof}

\begin{corollary}[Wulff 射线的精确预言机容量律与恢复位率常数]\label{cor:xi-time-part9x-serrin-ray-exact-capacity}
在定理 \ref{thm:xi-time-part9x-serrin-ray-fibonacci-semirings} 的设定下，固定分辨率 \(m\ge 1\)，并只观察尺度窗口
\[
0\le \rho<2^m.
\]
定义残余类读出
\[
r_m(\rho W_K):=[\rho]_{F_{m+2}}
\in
\ZZ/(F_{m+2}\ZZ),
\]
以及对应纤维基数
\[
d_m^{\mathrm{ray}}(r)
:=
\#\{\,0\le \rho<2^m:\ [\rho]_{F_{m+2}}=r\,\}.
\]
再记
\[
q_m:=\left\lfloor\frac{2^m}{F_{m+2}}\right\rfloor,
\qquad
a_m:=2^m-q_mF_{m+2}.
\]
则对每个 \(r\in \ZZ/(F_{m+2}\ZZ)\) 都有
\[
d_m^{\mathrm{ray}}(r)\in\{q_m,q_m+1\},
\]
且恰有 \(a_m\) 条纤维取值 \(q_m+1\)，其余 \(F_{m+2}-a_m\) 条纤维取值 \(q_m\)。

定义该射线上的 dyadic 预言机容量
\[
C_m^{\mathrm{ray}}(B)
:=
\sum_{r\in \ZZ/(F_{m+2}\ZZ)}
\min\bigl(d_m^{\mathrm{ray}}(r),2^B\bigr),
\qquad
B\ge 0.
\]
则有严格闭式
\[
C_m^{\mathrm{ray}}(B)
=
a_m\min(q_m+1,2^B)
+
\bigl(F_{m+2}-a_m\bigr)\min(q_m,2^B).
\]
相应最优成功率为
\[
\mathrm{Succ}_m^{\mathrm{ray}}(B)
=
\frac{C_m^{\mathrm{ray}}(B)}{2^m}.
\]
若要求在窗口 \(0\le \rho<2^m\) 内精确恢复全部尺度，则所需最小附加比特数满足
\[
B_m^{\mathrm{exact}}
=
\left\lceil \log_2(q_m+1)\right\rceil.
\]
并且存在极限
\[
\lim_{m\to\infty}\frac{B_m^{\mathrm{exact}}}{m}
=
\log_2\!\Bigl(\frac{2}{\varphi}\Bigr).
\]
\end{corollary}
\begin{proof}
将 \(2^m\) 作 Euclid 分解
\[
2^m=q_mF_{m+2}+a_m,
\qquad
0\le a_m<F_{m+2}.
\]
于是模 \(F_{m+2}\) 的每个剩余类在区间 \([0,2^m)\) 中出现 \(q_m\) 次或 \(q_m+1\) 次；并且恰有 \(a_m\) 个剩余类多出一次。这就得到纤维直方图的二值结构。

将该直方图代入与定理 \ref{thm:xi-capacity-tail-histogram-moment-complete-statistic} 相同的容量泛函
\[
C(B)=\sum_x\min(d(x),2^B),
\]
便立得所示闭式。

若要求在给定残余类之外再加 \(B\) 比特以恢复窗口内全部尺度，则每条纤维内部至多可用 \(2^B\) 个不同附加标签。故若
\[
2^B<q_m+1,
\]
则某条最大纤维中必有两个尺度无法区分，精确恢复不可能。反之，只要
\[
2^B\ge q_m+1,
\]
就可在每条纤维内部赋予互异附加标签，从而实现精确恢复。于是
\[
B_m^{\mathrm{exact}}
=
\left\lceil \log_2(q_m+1)\right\rceil.
\]

最后，由 Binet 公式
\[
F_{m+2}\sim \frac{\varphi^{m+2}}{\sqrt5}
\]
得到
\[
q_m
\sim
\frac{2^m}{F_{m+2}}
\sim
\frac{\sqrt5}{\varphi^2}
\Bigl(\frac{2}{\varphi}\Bigr)^m.
\]
故
\[
\log_2(q_m+1)
=
m\log_2\!\Bigl(\frac{2}{\varphi}\Bigr)+O(1),
\]
从而
\[
\frac{B_m^{\mathrm{exact}}}{m}\longrightarrow \log_2\!\Bigl(\frac{2}{\varphi}\Bigr).
\]
证毕。
\end{proof}

\begin{theorem}[局部化 Wulff Minkowski 射线的圆因子--素数账本分解]\label{thm:xi-time-part9x-serrin-ray-localized-ledger}
取有限素数集 \(S\subset\mathbb P\)，并定义同心局部化尺度半群
\[
\mathfrak W_{K,S}^{+}
:=
\{\rho W_K:\ \rho\in \ZZ[S^{-1}]_{\ge 0}\},
\]
其运算取为 Minkowski 和 \(\boxplus\)。记
\[
\mathfrak W_{K,S}^{\mathrm{gp}}
\]
为其 Grothendieck 群完成。则存在规范群同构
\[
\mathfrak W_{K,S}^{\mathrm{gp}}\cong \ZZ[S^{-1}].
\]
因此其 Pontryagin 对偶满足短正合列
\[
0\longrightarrow \prod_{p\in S}\ZZ_p
\longrightarrow
\widehat{\mathfrak W_{K,S}^{\mathrm{gp}}}
\longrightarrow
\TT
\longrightarrow 0.
\]
并且
\[
\cdim_{\mathrm{fr}}\bigl(\mathfrak W_{K,S}^{\mathrm{gp}}\bigr)=1,
\]
\[
p\in S
\iff
\text{核 }\prod_{q\in S}\ZZ_q\text{ 的 \(p\)-主因子非平凡}.
\]
\end{theorem}
\begin{proof}
对同心 Wulff 形有
\[
\rho W_K\boxplus \sigma W_K=(\rho+\sigma)W_K.
\]
因此映射
\[
\iota_S:\mathfrak W_{K,S}^{+}\longrightarrow \ZZ[S^{-1}]_{\ge 0},
\qquad
\iota_S(\rho W_K)=\rho
\]
是加法半群同构，进而其 Grothendieck 群完成满足
\[
\mathfrak W_{K,S}^{\mathrm{gp}}\cong \ZZ[S^{-1}].
\]

再由定理 \ref{thm:killo-localization-circle-dimension-prime-ledger-splitting} 作用于 \(G_S=\ZZ[S^{-1}]\)，立得 Pontryagin 对偶短正合列
\[
0\to \prod_{p\in S}\ZZ_p\to \widehat{G_S}\to \TT\to 0,
\]
以及
\[
\cdim_{\mathrm{fr}}(G_S)=1,
\qquad
p\in S
\iff
\text{核的 \(p\)-主因子非平凡}.
\]
经由 \(\mathfrak W_{K,S}^{\mathrm{gp}}\cong G_S\) 传回即得结论。证毕。
\end{proof}

\begin{proposition}[window-\texorpdfstring{$6$}{6} boundary 三点集的算术钉扎与非陪集性]\label{prop:xi-time-part9x-window6-boundary-arithmetic-pinning}
在定理 \ref{thm:xi-window6-cyclic-kernel-bc3-root-datum} 的分裂
\[
X_6=X_6^{\mathrm{cyc}}\sqcup X_6^{\mathrm{bdry}},
\qquad
|X_6|=21,\quad |X_6^{\mathrm{cyc}}|=18,\quad |X_6^{\mathrm{bdry}}|=3
\]
下，boundary 扇区满足
\[
X_6^{\mathrm{bdry}}
=
\{\texttt{100001},\texttt{100101},\texttt{101001}\}.
\]
通过定理 \ref{thm:fold-congruence-mod-semirings} 的同构
\[
\Omega_6/\!\equiv_6\cong \ZZ/(F_8\ZZ)=\ZZ/(21\ZZ),
\]
该三点集精确对应于剩余类集合
\[
\mathcal B_6=\{14,17,19\}\subset \ZZ/(21\ZZ).
\]
此外：
\[
\mathcal B_6\ \text{不是}\ \ZZ/(21\ZZ)\ \text{中的加法陪集};
\]
\[
\mathcal B_6\ \text{也不是单位群}\ (\ZZ/(21\ZZ))^\times\ \text{中任何乘法子群的陪集}.
\]
并且，对每个 \(w\in X_6^{\mathrm{bdry}}\)，其 dyadic 二层预像都是一个两点纤维，且两点差值恒为
\[
34=F_9.
\]
\end{proposition}
\begin{proof}
对稳定词 \(w=(w_1,\dots,w_6)\in X_6\)，其值映射为
\[
V_6(w):=\sum_{k=1}^{6}w_kF_{k+1}.
\]
故
\[
V_6(\texttt{100001})=F_2+F_7=1+13=14,
\]
\[
V_6(\texttt{101001})=F_2+F_4+F_7=1+3+13=17,
\]
\[
V_6(\texttt{100101})=F_2+F_5+F_7=1+5+13=19.
\]
于是
\[
\mathcal B_6=\{14,17,19\}.
\]

若 \(\mathcal B_6\) 是 \(\ZZ/(21\ZZ)\) 中的加法陪集，则因其基数为 \(3\)，必为唯一三阶加法子群
\[
\{0,7,14\}
\]
的某个陪集，因而任意两点差都应为 \(7\) 的倍数。但
\[
17-14=3,\qquad 19-17=2,\qquad 19-14=5,
\]
矛盾。故 \(\mathcal B_6\) 不是加法陪集。

另一方面，任何 \((\ZZ/(21\ZZ))^\times\) 中乘法子群的陪集都只由单位组成；但
\[
\gcd(14,21)=7\neq 1,
\]
故 \(14\notin(\ZZ/(21\ZZ))^\times\)。因此 \(\mathcal B_6\) 不可能是单位群中任何乘法子群的陪集。

最后，推论 \ref{cor:fold-bin-boundary-center-z2} 已给出：对每个
\[
w\in X_6^{\mathrm{bdry}},
\]
纤维 \((\Fold_6^{\mathrm{bin}})^{-1}(w)\) 都是二点集，并且两点差值固定为
\[
34=F_9.
\]
证毕。
\end{proof}

\begin{theorem}[boundary 扇区的 Weyl 完成阈值]\label{thm:xi-time-part9x-window6-boundary-weyl-threshold}
\leanverified{paper\_xi\_time\_part9x\_window6\_boundary\_weyl\_threshold}
称集合 \(Y\subset X_6\) 为 boundary 扇区的一个 Weyl 完成，若
\[
X_6^{\mathrm{bdry}}\subset Y
\]
且 \(Y\cap X_6^{\mathrm{cyc}}\) 在定理 \ref{thm:xi-time-part9r-window6-hamming-complete-weyl-orbit-invariant} 的 \(B_3/C_3\) 根字典下对 Weyl 群作用不变。则必有
\[
|Y|\in\{3,9,15,21\}.
\]
特别地，任何非平凡 Weyl 完成都满足
\[
|Y|\ge 9.
\]
\end{theorem}
\begin{proof}
由定理 \ref{thm:xi-time-part9r-window6-hamming-complete-weyl-orbit-invariant}，\(X_6^{\mathrm{cyc}}\) 在 \(B_3/C_3\) 两张字典下都恰分裂为两条 Weyl 轨道，其基数分别为
\[
6,\qquad 12.
\]
因此 \(X_6^{\mathrm{cyc}}\) 的 Weyl 不变子集只能是
\[
\varnothing,
\qquad
\text{六元轨道},
\qquad
\text{十二元轨道},
\qquad
X_6^{\mathrm{cyc}},
\]
其基数依次为
\[
0,\ 6,\ 12,\ 18.
\]

若 \(Y\) 含全部 boundary 三点，而 \(Y\cap X_6^{\mathrm{cyc}}\) Weyl 不变，则
\[
|Y|
=
|X_6^{\mathrm{bdry}}|+|Y\cap X_6^{\mathrm{cyc}}|
=
3+|Y\cap X_6^{\mathrm{cyc}}|
\in
\{3,9,15,21\}.
\]
若要求 completion 非平凡，即
\[
Y\cap X_6^{\mathrm{cyc}}\neq\varnothing,
\]
则最小可能值对应于加入六元轨道，故
\[
|Y|\ge 3+6=9.
\]
证毕。
\end{proof}

\begin{theorem}[六倍窗上 Wulff 射线算术骨架的半阶零点块]\label{thm:xi-time-part9x-serrin-ray-arithmetic-zero-block}
记
\[
\mathcal Z_m\subset \ZZ/(F_{m+2}\ZZ)
\]
为附录 \ref{app:fold-multiplicity} 的谱零点集合。通过定理
\ref{thm:xi-time-part9x-serrin-ray-fibonacci-semirings} 的有限分辨率同构，把它传送为 Wulff 射线算术骨架上的零点类集合
\[
\mathcal Z_{m,K}^{\mathrm{arith}}
:=
\bigl\{[\rho W_K]_{\equiv_m^{(K)}}:\ [\rho]_{F_{m+2}}\in \mathcal Z_m\bigr\}
\subset
\mathfrak W_K^{(0)}/\!\equiv_m^{(K)}.
\]
若
\[
m+2\equiv 0\pmod 6,
\qquad
d=\frac{m+2}{2},
\]
则
\[
\bigl|\mathcal Z_{m,K}^{\mathrm{arith}}\bigr|
\ge
F_d
=
F_{(m+2)/2}.
\]
特别地，在最小三轴窗
\[
m=58
\]
处，有精确值
\[
\bigl|\mathcal Z_{58,K}^{\mathrm{arith}}\bigr|
=
839415
\ge
F_{30}=832040.
\]
\end{theorem}
\begin{proof}
由定理 \ref{thm:xi-time-part9x-serrin-ray-fibonacci-semirings}，集合
\[
\mathfrak W_K^{(0)}/\!\equiv_m^{(K)}
\]
与
\[
\ZZ/(F_{m+2}\ZZ)
\]
之间存在规范双射，故
\[
\bigl|\mathcal Z_{m,K}^{\mathrm{arith}}\bigr|
=
|\mathcal Z_m|.
\]

若 \(m+2\equiv 0\pmod 6\)，则由推论 \ref{cor:fold-zero-half-index-multiple6}，
\[
|\mathcal Z_m|
\ge
F_{(m+2)/2}.
\]
于是
\[
\bigl|\mathcal Z_{m,K}^{\mathrm{arith}}\bigr|
\ge
F_{(m+2)/2}.
\]

当 \(m=58\) 时，注 \ref{rem:fold-zero-lb-m58} 给出
\[
|\mathcal Z_{58}|=839415.
\]
故同样有
\[
\bigl|\mathcal Z_{58,K}^{\mathrm{arith}}\bigr|=839415.
\]
证毕。
\end{proof}

\begin{conjecture}[单参数各向异性 Serrin 观测的 Fibonacci 同余普适性]\label{conj:xi-time-part9x-serrin-ray-fibonacci-universality}
设
\[
q_m:\mathfrak W_K^{(0)}\twoheadrightarrow Q_m
\qquad(m\ge 1)
\]
是一族附着于整数尺度 Wulff 射线的有限分辨率观测，并满足：
\begin{enumerate}
  \item \(Q_m\) 为有限交换半环，且
  \[
  |Q_m|=F_{m+2};
  \]
  \item \(q_m\) 精确实现 Minkowski 和与尺度复合，即
  \[
  q_m(x\boxplus y)=q_m(x)+q_m(y),
  \qquad
  q_m(x\boxtimes y)=q_m(x)\,q_m(y);
  \]
  \item 这些有限商在加法 Grothendieck 包络上相容，并诱导与
  \[
  \{F_{m+2}\ZZ:\ m\ge 1\}
  \]
  相同的 Fibonacci 同余拓扑。
\end{enumerate}
则应存在唯一一族半环同构
\[
Q_m\cong \ZZ/(F_{m+2}\ZZ)
\qquad(m\ge 1),
\]
使之与全部相容结构交换。等价地，任何满足上述三条公理的单参数各向异性 Serrin 有限分辨率观测，其完备化都应退化为同一个 Fibonacci 同余完备化
\[
\widehat{\ZZ}_F\cong \widehat{\ZZ}.
\]
若该猜想成立，则
\[
\log_2\!\Bigl(\frac{2}{\varphi}\Bigr)
\]
将成为单参数可实现几何的普适精确恢复指数，而定理
\ref{thm:xi-time-part9x-serrin-ray-arithmetic-zero-block} 的半阶零点块则给出其在六倍窗上的普适算术盲区指纹。
\end{conjecture}
\begin{proof}[证据]
定理 \ref{thm:xi-time-part9x-serrin-ray-fibonacci-semirings} 已经给出一条显式模型：整数尺度 Wulff 射线在每个分辨率 \(m\) 上都产生模 \(F_{m+2}\) 半环。推论 \ref{cor:xi-time-part9x-serrin-ray-exact-capacity} 说明，一旦这一算术骨架被固定，精确恢复阈值便被强制为
\[
\log_2(2/\varphi).
\]
定理 \ref{thm:xi-time-part9x-serrin-ray-localized-ledger} 进一步说明，该骨架的局部化群完成在对偶端必裂解为一个圆因子与完全可恢复的 \(p\)-adic 账本核；而定理 \ref{thm:xi-time-part9x-serrin-ray-arithmetic-zero-block} 则表明，在六倍窗子序列上同一骨架还携带一个确定性的半阶零点类块。

因此，上述猜想所断言的是：对单参数各向异性 Serrin 观测而言，一旦有限分辨率对象数、两种运算以及诱导拓扑同时被固定，就不再存在第二种与之不同的算术骨架。换言之，Fibonacci 模数族不只是一条可实现模型，而应是这一单参数几何在有限分辨率端的唯一刚性终型。证毕。
\end{proof}

\noindent
上述链条表明，未归一化 Wulff 射线在有限分辨率上并不退化为平凡的单参数残余；相反，它携带一条精确可算的 Fibonacci 同余骨架。该骨架在局部化群完成后裂解为一个圆因子与一个可完全反演的 \(p\)-adic 账本核，在 window-\(6\) 上把 boundary 三点集精确钉扎为 \(\ZZ/(21\ZZ)\) 中一个既非加法陪集亦非乘法陪集的三点残余，并在六倍窗子序列上承载一个半阶量级的算术零点块。因而，对单参数可实现几何而言，有限分辨率容量律、局部化对偶核、boundary 阈值与六倍窗盲区不再是彼此分离的现象，而应被视为同一条 Fibonacci 同余主线的不同投影。

\endinput
